\documentclass[11pt,a4paper]{article}
\usepackage[T1]{fontenc}
\usepackage[utf8]{inputenc}
\usepackage[brazilian]{babel}
\usepackage[margin=1in]{geometry}
\usepackage{setspace}
\usepackage{enumitem}
\usepackage{array}
\usepackage{hyperref}
\usepackage{xurl}
\usepackage{parskip}
\usepackage{longtable}
\usepackage{fvextra}
\usepackage{csquotes}
\usepackage{tikz}
\usetikzlibrary{arrows.meta, positioning}
\usepackage[
  backend=biber,
  style=abnt,
  sorting=nty,
  maxcitenames=2,
  maxbibnames=99,
  uniquelist=false,
  repeatfields=false,
  noslsn=true
]{biblatex}
\addbibresource{chatbot-rag.bib}

\setstretch{1.15}
\setlength{\parindent}{0pt}
\emergencystretch=2em
\setlist[enumerate]{leftmargin=2em}
\hypersetup{colorlinks=true,linkcolor=black,urlcolor=blue,citecolor=blue}

\begin{document}

\begin{center}
{\large Instituto Brasileiro de Ensino, Desenvolvimento e Pesquisa}\par
Mestrado em Administração Pública -- Ciência de Dados e Inteligência Artificial\par
Turma: 2005.2\par
Disciplina: Inteligência Artificial Generativa\par
Professor: Marcelo Pita\par
Alunos: Analécia Rorato; Fabricio Fernandes Santana; Paulo Ceser Siqueira Neves

\vspace{1.5em}
{\LARGE \textbf{Chatbot legislativo com RAG para respostas auditáveis sobre discursos da 56\textordfeminine{} Legislatura do Senado Federal}}\par

\vspace{0.75em}
Repositório: \url{https://github.com/fabriciosantana/mcdia/tree/main/05-iag/4-project}
\end{center}

\section*{\centering Resumo}

Este trabalho investiga a aplicação de \textit{retrieval-augmented generation} (RAG) no desenvolvimento de um \textit{chatbot} legislativo para consulta a discursos da 56\textordfeminine{} Legislatura do Senado Federal (2019--2023). Parte-se da dificuldade de recuperar informações parlamentares por meio de linguagem natural com precisão factual, contextualização e rastreabilidade. A metodologia consistiu na construção de um \textit{pipeline} com preparação de dados, segmentação textual, enriquecimento por metadados, geração de \textit{embeddings}, indexação vetorial e fluxo de consulta baseado em LLM. A solução foi implementada em ambiente reprodutível, com \textit{scripts} próprios para ingestão e avaliação, e validada por três estudos de caso manuais, três rodadas automatizadas com rubrica e LLM como juiz, além de revisão humana amostral das respostas automatizadas. Os resultados indicam bom desempenho em perguntas factuais, temáticas e de autoria bem delimitada, com médias entre 9,15 e 9,55 nas rodadas automatizadas. Observou-se maior instabilidade em consultas que exigem precisão numérica, comparação entre autores, controle de escopo ou síntese multifocal. Conclui-se que a abordagem é viável para recuperação auditável de discursos legislativos, desde que acompanhada de governança de fontes, avaliação sistemática e supervisão humana.

\textbf{Palavras-chave:} RAG; inteligência artificial; \textit{chatbot} legislativo; recuperação de informação; Senado Federal; LLM.

\section{Introdução}

\subsection{Contextualização}

O avanço dos grandes modelos de linguagem (\textit{large language models} -- LLMs) ampliou significativamente o interesse por soluções capazes de responder perguntas em linguagem natural \parencite{minaee2024llmSurvey}. No setor público, esse interesse é particularmente relevante diante de acervos extensos, dinâmicos e heterogêneos, tais como legislação, pareceres, notas técnicas, discursos parlamentares e documentos administrativos \parencite{ipu2024aiParliaments}. Nesses ambientes, o desafio não se restringe ao armazenamento da informação, mas abrange, sobretudo, sua recuperação qualificada, contextualizada e verificável \parencite{wfd2023guidelinesAiParliaments}.

Apesar de seu potencial, o uso direto de LLMs apresenta limitações amplamente discutidas na literatura, entre as quais se destacam a alucinação, a dependência do conhecimento incorporado aos dados de treinamento e a dificuldade de identificar as evidências utilizadas na produção das respostas \parencite{gao2023ragSurvey}. Tais limitações assumem maior gravidade em contextos institucionais nos quais se exigem transparência, verificabilidade, controle das fontes e responsabilização pelos conteúdos produzidos \parencite{wfd2023guidelinesAiParliaments}.

Nesse cenário, a abordagem \textit{retrieval-augmented generation} (RAG) apresenta-se como alternativa tecnicamente promissora, pois, em vez de depender exclusivamente do conhecimento internalizado no modelo, sistemas RAG recuperam trechos relevantes de uma base documental externa e os incorporam ao contexto da geração \parencite{gao2023ragSurvey}. Com isso, a abordagem pode elevar a precisão factual e a atualidade das respostas, ao mesmo tempo em que favorece a explicitação das fontes que as sustentam \parencite{gao2023ragSurvey}.

No contexto do Parlamento brasileiro, a organização do conhecimento legislativo e a ampliação do acesso qualificado à informação constituem desafios relevantes de pesquisa aplicada. Como observam \textcite{martelloteViola2025organizacaoConhecimento}, o emprego de inteligência artificial nesse domínio deve ser acompanhado de atenção à estrutura informacional, à mediação documental e à explicabilidade dos resultados.

Com base nessas questões, o presente trabalho propõe o desenvolvimento de um \textit{chatbot} legislativo baseado em RAG, voltado à consulta de discursos da 56\textordfeminine{} Legislatura do Senado Federal.

\subsection{Motivação}

A administração pública contemporânea opera sob crescente pressão por transparência, rastreabilidade, eficiência e capacidade de resposta \parencite{un2015responsiveGovernance}. No ambiente parlamentar, a consulta a discursos e pronunciamentos mostra-se relevante para diversas finalidades, tais como análise temática, comparação de posicionamentos, apoio técnico a gabinetes, elaboração de relatórios, produção de subsídios e atendimento ao cidadão \parencite{skubicFiser2024discourseReview}. Estudos da Consultoria Legislativa do Senado mostram, além disso, que os pronunciamentos em plenário entre 2007 e 2024 apresentam transformações relevantes em volume, extensão, interatividade e complexidade textual, reforçando a pertinência de tratá-los como corpus dinâmico de pesquisa institucional \parencite{blanco2025plenarioPalanqueEstudio,blanco2026letrasGraudas}. Entretanto, a simples disponibilização pública desses documentos não garante, por si só, acesso eficiente à informação que eles contêm \parencite{geunis2023parliamentaryMonitoring}.

Em grandes acervos discursivos, mecanismos tradicionais de busca, predominantemente baseados em correspondência lexical, tendem a ser insuficientes quando o usuário formula perguntas complexas, comparativas, interpretativas ou orientadas à síntese \parencite{gupta2024comprehensiveRag}. Além disso, em bases documentais extensas, a recuperação da informação frequentemente depende de conhecimento prévio sobre vocabulário, datas, autores ou sobre a própria organização do acervo, o que dificulta o acesso substantivo à informação mesmo quando ela está formalmente disponível \parencite{bandeiraBernardes2016information,geunis2023parliamentaryMonitoring}. Como consequência, a informação permanece formalmente acessível, mas substantivamente difícil de explorar \parencite{bandeiraBernardes2016information,geunis2023parliamentaryMonitoring}.

Paralelamente, soluções baseadas exclusivamente em LLMs generalistas podem produzir respostas plausíveis e linguisticamente bem estruturadas, porém desprovidas de fundamentação verificável \parencite{anhHoang2025hallucinations,gao2023ragSurvey}. Em domínios institucionais, tal comportamento é particularmente problemático, pois compromete a confiança, a responsabilização e o uso controlado da tecnologia \parencite{deAlmeidaSantos2025aiGovernance,wfd2023guidelinesAiParliaments}. A literatura de RAG indica que parte desses riscos pode ser mitigada ao ancorar a geração em documentos externos, controláveis e atualizáveis, desde que isso seja combinado a práticas de avaliação, versionamento, rastreabilidade e monitoramento contínuo \parencite{gao2023ragSurvey,saadFalconEtAl2024ares,esEtAl2023ragas}.

Dessa forma, este trabalho é motivado pela necessidade de investigar uma solução tecnicamente viável e institucionalmente adequada para ampliar o acesso qualificado a discursos parlamentares, combinando busca semântica, geração textual e verificação documental.

\subsection{Problema de pesquisa}

Diante desse quadro, formula-se o seguinte problema de pesquisa: como obter, por meio de consultas em linguagem natural aos discursos da 56\textordfeminine{} Legislatura do Senado Federal, respostas úteis, fundamentadas em evidências recuperadas e compatíveis com exigências de verificação institucional?

Esse problema de pesquisa decorre da constatação de que o acesso documental, embora formalmente assegurado, ainda enfrenta limitações relevantes quando se exigem recuperação semântica, síntese orientada ao usuário e indicação clara das fontes utilizadas na resposta \parencite{bandeiraBernardes2016information,geunis2023parliamentaryMonitoring,gao2023ragSurvey}.

\subsection{Objetivos}

O objetivo geral deste trabalho é desenvolver e demonstrar uma prova de conceito reprodutível de \textit{chatbot} legislativo baseada em RAG para consulta a discursos parlamentares da 56\textordfeminine{} Legislatura do Senado Federal, com foco na produção de respostas fundamentadas em evidências recuperadas e compatíveis com exigências de verificação institucional.

Como objetivos específicos, definem-se:
\begin{enumerate}
  \item estruturar uma base de conhecimento consultável a partir dos discursos parlamentares;
  \item possibilitar a formulação de perguntas em linguagem natural com respostas fundamentadas em trechos recuperados;
  \item avaliar empiricamente a qualidade inicial da recuperação e das respostas geradas;
  \item documentar uma arquitetura reprodutível, modular e passível de evolução para cenários institucionais.
\end{enumerate}

Em razão de sua natureza construtiva, o estudo não formula hipóteses experimentais no sentido clássico. A avaliação é orientada por critérios de verificação associados à construção e demonstração de um artefato computacional: base documental indexável, recuperação de evidências, geração fundamentada, rastreabilidade das fontes, automação do fluxo e reprodutibilidade. Essa orientação aproxima o trabalho da \textit{Design Science Research}, na qual a produção e a avaliação de artefatos constituem forma legítima de investigação em sistemas de informação \parencite{hevnerEtAl2004designScience,peffersEtAl2007dsrm}.

\subsection{Justificativa}

A relevância deste estudo decorre de sua contribuição em dois planos complementares. No plano prático, a pesquisa investiga uma forma de ampliar o acesso qualificado a informações parlamentares públicas, com potencial para fortalecer transparência, memória institucional e apoio a atividades de gabinete, pesquisa e controle social \parencite{ipu2024aiParliaments,geunis2023parliamentaryMonitoring}. No plano metodológico e aplicado, o trabalho examina a aplicação de técnicas contemporâneas de recuperação e geração em um domínio documental complexo, marcado por exigências elevadas de explicabilidade, confiabilidade e vinculação às fontes \parencite{martelloteViola2025organizacaoConhecimento,wfd2023guidelinesAiParliaments}.

Além disso, o estudo se justifica pela aderência do problema a desafios concretos da administração pública digital. Em vez de tratar a inteligência artificial como mecanismo autônomo de produção de respostas, propõe-se uma abordagem em que a geração textual é condicionada por fontes documentais públicas, controláveis e auditáveis. Essa orientação é coerente com demandas institucionais por uso responsável da IA em contextos sensíveis \parencite{deAlmeidaSantos2025aiGovernance,wfd2023guidelinesAiParliaments}.

\section{Fundamentação teórica}

\subsection{Arquiteturas de LLM relacionadas ao problema}

Os componentes mais diretamente relacionados a este trabalho são, de um lado, os LLMs generativos baseados em transformers \parencite{vaswani2017attention} e, de outro, os modelos de embeddings utilizados para recuperação semântica \parencite{nie2024textEmbeddingSurvey}. Enquanto o LLM atua como componente de síntese textual, os embeddings convertem documentos e consultas em vetores, permitindo o cálculo de similaridade semântica em bancos vetoriais \parencite{gao2023ragSurvey}.

Em sistemas RAG, esses componentes operam de maneira complementar: o modelo de embedding sustenta a recuperação, enquanto o modelo generativo formula a resposta a partir do contexto recuperado \parencite{lewisEtAl2020rag,gao2023ragSurvey}. Essa separação favorece uma arquitetura modular, a especialização de componentes e a evolução incremental da solução \parencite{gao2024modularRag}.

Segundo \textcite{gao2023ragSurvey}, sistemas RAG podem ser classificados em abordagens ingênuas, avançadas e modulares. Para o problema investigado, a arquitetura modular mostra-se adequada por permitir substituir, de forma independente, extratores documentais, estratégias de segmentação, modelos de embedding, bancos vetoriais, reranqueadores e modelos gerativos. Em contextos institucionais, essa característica é desejável porque infraestrutura, custos, segurança e políticas internas podem exigir alterações graduais da pilha tecnológica.

Essa modularidade, contudo, não deve ser interpretada como garantia automática de qualidade. A literatura apresenta o RAG como estratégia para reduzir problemas de desatualização, baixa verificabilidade e alucinação, mas não como solução suficiente para eliminá-los. O desempenho final continua dependente da qualidade documental, da segmentação, da recuperação, do desenho do prompt, da seleção de contexto e dos mecanismos de avaliação adotados \parencite{gao2023ragSurvey,gupta2024comprehensiveRag,saadFalconEtAl2024ares}. Essa leitura é importante em aplicações públicas, nas quais uma resposta plausível, mas mal fundamentada, pode comprometer confiança institucional.

\subsection{Técnicas relacionadas}

As principais técnicas relacionadas ao \textit{pipeline} deste projeto incluem extração documental, \textit{chunking}, geração de \textit{embeddings} semânticos, recuperação vetorial, seleção contextual e \textit{prompting} orientado a evidências, podendo também envolver etapas de reranqueamento em arquiteturas RAG \parencite{gao2023ragSurvey,maharjan2026chunkingReranking,wangEtAl2025segmentation}. Em conjunto, essas técnicas convertem o acervo em unidades recuperáveis, identificam trechos semanticamente relevantes e delimitam o contexto encaminhado ao modelo generativo. O \textit{prompting} orientado a evidências complementa esse fluxo ao instruir o modelo a privilegiar o contexto recuperado, evitar inferências especulativas e explicitar referências verificáveis.

Dentre essas técnicas, o \textit{chunking} ocupa posição central. Estudos sobre segmentação documental em RAG indicam que escolhas inadequadas de divisão podem comprometer a recuperação semântica, seja por fragmentarem indevidamente o contexto, seja por gerarem blocos excessivamente extensos, com perda de precisão na busca \parencite{wangEtAl2025segmentation,bhatEtAl2025chunkSize}. No presente projeto, a segmentação documental exigiu sucessivos ajustes empíricos para equilibrar granularidade semântica, preservação de contexto e restrições operacionais da base vetorial e do processo de ingestão. Assim, a experiência prática corrobora a literatura ao indicar que o \textit{chunking} não constitui etapa meramente operacional, mas um dos principais determinantes do desempenho de um sistema RAG.

\textit{Surveys} e \textit{frameworks} também indicam que, em ambientes reais, a qualidade de sistemas RAG depende de mecanismos de avaliação, rastreabilidade e comparação entre configurações, não apenas do \textit{pipeline} básico de recuperação e geração \parencite{gao2023ragSurvey,saadFalconEtAl2024ares,esEtAl2023ragas}. Essa perspectiva orientou a criação de \textit{scripts} próprios de preparação, importação, avaliação e documentação da configuração experimental.

No campo da avaliação, estratégias baseadas em LLM como juiz têm sido utilizadas para escalar a análise de respostas geradas, sobretudo quando combinadas com rubricas explícitas \parencite{liuEtAl2023geval,liuEtAl2025judgeRag}. Ainda assim, o julgamento automatizado pode variar conforme o modelo avaliador adotado, a rubrica utilizada e o grau de abertura da tarefa; por isso, em aplicações RAG, tende a ser mais robusto quando tratado como instrumento auxiliar, combinado a evidências qualitativas, validação humana e análise das fontes recuperadas \parencite{liuEtAl2025judgeRag,liuEtAl2023geval,saadFalconEtAl2024ares,esEtAl2023ragas}. Benchmarks voltados ao domínio jurídico reforçam essa preocupação ao tratar a avaliação de RAG como problema que envolve não apenas qualidade textual da resposta, mas também recuperação de documentos pertinentes, fidelidade ao contexto e adequação da evidência utilizada \parencite{pipitoneAlami2024legalbenchRag}.

\subsection{Trabalhos e soluções relacionadas}

\textcite{gao2023ragSurvey} oferecem uma base conceitual abrangente para compreender o RAG como resposta aos problemas de desatualização, baixa verificabilidade e geração não rastreável em LLMs, além de orientar a decomposição do problema em etapas de recuperação, aumento de contexto e geração. No plano arquitetural, \textcite{gupta2024comprehensiveRag} consolidam o estado da arte em RAG ao discutir como mecanismos de recuperação e geração se combinam para mitigar limitações de LLMs e ao destacar desafios de precisão, fidelidade e coordenação entre componentes.

Os trabalhos de \textcite{minaee2024llmSurvey} e \textcite{vaswani2017attention} oferecem o pano de fundo técnico para o uso de LLMs generativos neste projeto: o primeiro sintetiza a evolução e as capacidades desses modelos em ampla gama de tarefas, enquanto o segundo estabelece a base arquitetural dos transformers. \textcite{nie2024textEmbeddingSurvey} complementam esse quadro ao tratar os \textit{text embeddings} como tecnologia fundacional para tarefas de busca semântica e recuperação de informação, enfatizando a relevância da integração entre LLMs e \textit{embeddings} em aplicações contemporâneas.

A literatura sobre alucinação em LLMs evidencia que esses modelos podem produzir saídas fluentemente redigidas, porém factualmente incorretas ou não verificáveis, o que reforça, em contextos públicos sensíveis, a necessidade de soluções explicáveis e auditáveis, com mecanismos de fundamentação documental, governança e controle institucional \parencite{anhHoang2025hallucinations,gao2023ragSurvey,deAlmeidaSantos2025aiGovernance}. Em domínios jurídicos, esse risco é particularmente relevante porque respostas plausíveis podem induzir interpretações normativas, citações inexistentes ou inferências de autoridade sem base documental suficiente \parencite{dahlEtAl2024largeLegalFictions}.

Em termos de contexto institucional, \textcite{un2015responsiveGovernance}, \textcite{ipu2024aiParliaments}, \textcite{wfd2023guidelinesAiParliaments} e \textcite{geunis2023parliamentaryMonitoring} indicam que parlamentos e administrações públicas operam sob pressões crescentes por transparência, responsividade, responsabilização e uso qualificado da informação. Nesse domínio, discursos e debates constituem insumos centrais para análise política e apoio técnico, mas os sítios parlamentares frequentemente funcionam mais como repositórios informacionais do que como instrumentos de exploração substantiva, tornando o acesso dependente de conhecimento prévio sobre estrutura e vocabulário \parencite{skubicFiser2024discourseReview,bandeiraBernardes2016information}.

No caso brasileiro, os Textos para Discussão da Consultoria Legislativa do Senado oferecem evidência empírica sobre a relevância e a complexidade do acervo de pronunciamentos. \textcite{blanco2025plenarioPalanqueEstudio} analisa discursos em plenário entre 2007 e 2024 e identifica mudanças de volume, extensão, tempo de fala e interatividade, associadas a transformações institucionais e tecnológicas. Em estudo complementar, \textcite{blanco2026letrasGraudas} examina leiturabilidade e traços morfossintáticos do mesmo domínio discursivo por meio de técnicas computacionais de análise textual. Esses trabalhos reforçam que os discursos do Senado constituem corpus relevante para métodos computacionais, embora se concentrem em análise descritiva e interpretativa, e não no desenvolvimento de uma camada interativa voltada à consulta, à recuperação semântica e à geração de respostas verificáveis.

No campo da governança pública de inteligência artificial, estudos e diretrizes institucionais enfatizam responsabilidade, transparência, monitoramento, gestão de riscos e alinhamento a finalidades públicas \parencite{deAlmeidaSantos2025aiGovernance,ipu2024aiParliaments,wfd2023guidelinesAiParliaments}. Para aplicações RAG em acervos legislativos, esse conjunto de referências reforça a importância de bases documentais controladas, metadados preservados, respostas verificáveis e avaliação contínua.

No domínio político, \textcite{khaliq2024ragar} demonstram que arquiteturas RAG podem apoiar tarefas de fact-checking político com raciocínio fundamentado em evidências externas, articulando recuperação rigorosa e geração apoiada em fontes verificáveis. Ainda que o objetivo deste trabalho não seja a verificação factual multimodal, a pesquisa desses autores é relevante por demonstrar a aderência do domínio político-parlamentar a abordagens baseadas em recuperação e geração com evidências.

Por fim, \textcite{martelloteViola2025organizacaoConhecimento} aproxima a discussão do contexto brasileiro ao tratar da inteligência artificial na organização do conhecimento legislativo. Sua contribuição reforça a pertinência de soluções explicáveis, documentalmente ancoradas e informacionalmente estruturadas em ambientes parlamentares, oferecendo referencial específico para a aplicação proposta neste trabalho.

Trabalhos recentes tornam mais visível a aproximação entre RAG, LLMs, conhecimento legislativo e acesso a informação parlamentar. A Tabela~\ref{tab:trabalhos-relacionados} sintetiza os estudos mais próximos identificados na revisão e explicita sua relação com a presente proposta. Em conjunto, eles mostram que há aplicações relevantes em fact-checking político, redação legislativa, recuperação de legislação, interfaces conversacionais e debates parlamentares \parencite{khaliq2024ragar,chouhanGertz2024lexdrafter,rogiersEtAl2024kamerraad,matoshiEtAl2025parliamentRag,garrisonEtAl2024talkNdaa,colombo2024legisAi,colomboEtAl2025legisSearch,mosbachEtAl2025worldAvatarParliament,sarnikar2025legislativeTexts}. Ainda assim, diferem do presente trabalho por privilegiarem outros corpus, línguas, jurisdições, tarefas ou estratégias de avaliação.

\begin{scriptsize}
\begin{longtable}{@{}p{0.17\textwidth}p{0.22\textwidth}p{0.22\textwidth}p{0.28\textwidth}@{}}
\caption{Trabalhos relacionados mais próximos e posicionamento da proposta}
\label{tab:trabalhos-relacionados}\\
\hline
\textbf{Trabalho} & \textbf{Corpus ou domínio} & \textbf{Abordagem e avaliação} & \textbf{Relação com este artigo} \\
\hline
\endfirsthead
\hline
\textbf{Trabalho} & \textbf{Corpus ou domínio} & \textbf{Abordagem e avaliação} & \textbf{Relação com este artigo} \\
\hline
\endhead
\textcite{blanco2025plenarioPalanqueEstudio} e \textcite{blanco2026letrasGraudas} & Discursos em plenário do Senado Federal brasileiro. & Análises computacionais descritivas sobre volume, extensão, interatividade, leiturabilidade e complexidade textual. & Compartilham o domínio empírico brasileiro, mas não propõem sistema RAG, interface conversacional ou avaliação de respostas geradas. \\
\textcite{martelloteViola2025organizacaoConhecimento} & Dados e informações legislativas do Parlamento brasileiro. & Discussão aplicada sobre IA e organização do conhecimento legislativo. & Aproxima o problema ao contexto brasileiro de organização da informação, mas não avalia um chatbot RAG sobre discursos parlamentares. \\
\textcite{khaliq2024ragar} & Fact-checking político multimodal. & Arquitetura RAG para raciocínio apoiado em evidências externas. & Demonstra a utilidade de recuperação com evidências no domínio político, embora com foco em verificação factual, e não em consulta a acervo legislativo. \\
\textcite{chouhanGertz2024lexdrafter} & Documentos legislativos da União Europeia, com foco em definições terminológicas. & RAG para apoiar redação de artigos de definição; código disponibilizado pelos autores. & É diretamente legislativo e usa RAG, mas a tarefa principal é redação terminológica, não consulta auditável a discursos parlamentares. \\
\textcite{rogiersEtAl2024kamerraad} & Política nacional belga. & Sumarização hierárquica e interface conversacional para recuperação de informação política. & Aproxima-se pela interface conversacional em domínio parlamentar, mas difere por corpus, língua, desenho de sumarização e foco avaliativo. \\
\textcite{matoshiEtAl2025parliamentRag} & Acervos e negócios parlamentares suíços. & Projetos-piloto com modelagem de tópicos e RAG para melhorar acesso digital. & É uma das aproximações mais diretas ao problema de acesso parlamentar, mas em outra jurisdição e com foco em indexação temática institucional. \\
\textcite{garrisonEtAl2024talkNdaa} & Legislação congressional dos Estados Unidos, em especial o NDAA. & Avaliação de RAG para sumarização e consulta de legislação longa. & Contribui para comparação em legislação e avaliação de RAG, mas não trata discursos parlamentares nem o contexto brasileiro. \\
\textcite{colomboEtAl2025legisSearch} & Legislação italiana representada em grafo. & Busca legislativa com grafo, LLM, embeddings, comparação com BM25/TF-IDF e ablação. & É metodologicamente forte para recuperação legislativa, mas prioriza navegação em leis e relações normativas, não discursos e metadados parlamentares. \\
\textcite{mosbachEtAl2025worldAvatarParliament} & Debates parlamentares do Bundestag alemão. & Sistema RAG híbrido com grafo de conhecimento e recuperação vetorial. & Aproxima-se por tratar debates parlamentares, mas usa arquitetura híbrida com grafo e outro contexto linguístico-institucional. \\
\textcite{sarnikar2025legislativeTexts} & Textos legislativos longos. & Uso de LLMs para consulta e compreensão de legislação extensa. & Dialoga com a dificuldade de consultar textos legislativos longos, mas não enfatiza corpus de discursos, metadados parlamentares ou prova de conceito reprodutível. \\
\hline
\end{longtable}
\end{scriptsize}

\subsection{Síntese da revisão e lacuna de pesquisa}

A revisão realizada permite identificar três conjuntos de literatura complementares. O primeiro, voltado a RAG, LLMs, \textit{embeddings}, segmentação documental e avaliação automatizada, oferece a base técnica do trabalho \parencite{lewisEtAl2020rag,gao2023ragSurvey,minaee2024llmSurvey,nie2024textEmbeddingSurvey,saadFalconEtAl2024ares}. O segundo, dedicado a parlamentos, organização da informação legislativa e discurso parlamentar, delimita a relevância institucional do domínio \parencite{skubicFiser2024discourseReview,martelloteViola2025organizacaoConhecimento,ipu2024aiParliaments}. O terceiro, mais diretamente ligado ao caso brasileiro, aplica métodos computacionais ao corpus de discursos do Senado, mas com foco predominante em análise descritiva, linguística ou discursiva \parencite{blanco2025plenarioPalanqueEstudio,blanco2026letrasGraudas}. 

A lacuna remanescente, portanto, não está na ausência de estudos sobre RAG nem na ausência de estudos sobre discursos parlamentares, mas na articulação aplicada entre esses dois campos. Embora a literatura indique o potencial do RAG para elevar a precisão factual e a verificabilidade de sistemas baseados em LLM, esse potencial depende de decisões metodológicas situadas, e não apenas da adoção abstrata da arquitetura \parencite{gao2023ragSurvey,gupta2024comprehensiveRag,saadFalconEtAl2024ares}. Parte significativa da literatura técnica permanece concentrada em discussões conceituais amplas, cenários corporativos ou aplicações em domínios distintos do acervo legislativo brasileiro \parencite{gao2023ragSurvey,gupta2024comprehensiveRag}. Mesmo os trabalhos mais próximos em legislação, política e parlamentos enfatizam fact-checking, redação de normas, recuperação de legislação, grafos de conhecimento ou outras jurisdições \parencite{khaliq2024ragar,chouhanGertz2024lexdrafter,rogiersEtAl2024kamerraad,matoshiEtAl2025parliamentRag,garrisonEtAl2024talkNdaa,colomboEtAl2025legisSearch,mosbachEtAl2025worldAvatarParliament,sarnikar2025legislativeTexts}. Ainda que pesquisas já explorem computacionalmente discursos do Senado sob perspectivas quantitativas, discursivas e linguísticas, são menos frequentes estudos que articulem, em um mesmo desenho aplicado, corpus parlamentar nacional, preparação reprodutível de base documental, preservação de metadados legislativos, fluxo RAG operacional e avaliação empírica das respostas geradas \parencite{blanco2025plenarioPalanqueEstudio,blanco2026letrasGraudas,martelloteViola2025organizacaoConhecimento}.

Assim, identifica-se uma lacuna aplicada relacionada à construção, à documentação e à avaliação de soluções RAG voltadas especificamente à consulta de discursos parlamentares no contexto do Senado Federal. Este trabalho busca contribuir para seu preenchimento ao propor, implementar e avaliar uma prova de conceito reprodutível de \textit{chatbot} legislativo baseada em RAG, com foco em discursos da 56\textordfeminine{} Legislatura. A contribuição reside não apenas no protótipo, mas também na explicitação do \textit{pipeline}, no enriquecimento dos \textit{chunks} com metadados, na documentação dos artefatos de ingestão e avaliação e na combinação entre estudos de caso, rodadas automatizadas com rubrica e análise manual amostral.

\section{Caracterização do problema}

\subsection{Problema da administração pública}

O problema central abordado neste trabalho consiste na dificuldade de transformar grandes acervos de discursos legislativos em bases efetivamente consultáveis por meio de perguntas em linguagem natural, com respostas úteis, contextualizadas e verificáveis. Embora esses discursos sejam públicos, seu volume, diversidade temática, dispersão temporal e variação formal tornam a recuperação de argumentos, posicionamentos, temas recorrentes e referências históricas uma tarefa complexa quando realizada por meio de mecanismos tradicionais de busca \parencite{kavcic2024historicalViewer,blanco2025plenarioPalanqueEstudio,blanco2026letrasGraudas}.

Na prática, esse quadro impacta diferentes atores institucionais. Gabinetes parlamentares podem precisar recuperar rapidamente posicionamentos anteriores sobre determinado tema; equipes técnicas precisam localizar subsídios para notas, relatórios e pareceres; pesquisadores e jornalistas demandam acesso qualificado a discursos para fins analíticos; e cidadãos interessados em transparência e controle social podem se beneficiar de instrumentos que tornem o acervo mais acessível e inteligível \parencite{bandeiraBernardes2016information}. Na ausência de uma camada de recuperação semântica, o acesso ao conteúdo tende a ser fragmentado, lento e altamente dependente de conhecimento prévio sobre a estrutura documental e os termos empregados nos discursos \parencite{gao2023ragSurvey,martelloteViola2025organizacaoConhecimento,geunis2023parliamentaryMonitoring}.

Desse modo, o \textit{chatbot} legislativo com RAG proposto neste trabalho procura enfrentar essa lacuna ao acrescentar uma camada de mediação informacional que combina recuperação semântica, contextualização documental e geração textual ancorada em evidências.

\subsection{Partes interessadas}

As partes interessadas no projeto abrangem, em primeiro lugar, gabinetes parlamentares, consultorias legislativas e assessorias técnicas, que demandam acesso rápido e qualificado a discursos e posicionamentos anteriores. Também se incluem equipes de transparência, gestão da informação e tecnologia, para as quais são particularmente relevantes aspectos de governança, manutenção, rastreabilidade e evolução da solução.

Além disso, pesquisadores das áreas de ciência política, direito, administração pública e ciência da informação constituem público potencial do sistema, bem como cidadãos e organizações da sociedade civil interessados em controle social e acompanhamento da atividade parlamentar \parencite{martelloteViola2025organizacaoConhecimento,geunis2023parliamentaryMonitoring}.

\subsection{Critérios de sucesso}

Os critérios de sucesso definidos para este trabalho incluem os seguintes aspectos: a capacidade de indexar uma base representativa de discursos parlamentares; a recuperação de trechos relevantes para perguntas factuais e analíticas; a geração de respostas em português fundamentadas no contexto recuperado; a apresentação de referências e metadados associados aos discursos utilizados; a automatização das etapas de importação e avaliação; e a reprodutibilidade da solução.

Tais critérios dialogam com recomendações da literatura para sistemas RAG e avaliação automatizada, que enfatizam avaliação contínua, reprodutibilidade e rastreabilidade como propriedades desejáveis em soluções voltadas ao tratamento de informação institucionalmente relevante \parencite{gao2023ragSurvey,saadFalconEtAl2024ares,esEtAl2023ragas}.

\section{Proposta de solução}

\subsection{Enquadramento metodológico}

O trabalho adota uma orientação construtiva e aplicada, compatível com a \textit{Design Science Research} em sistemas de informação. Nessa perspectiva, a contribuição científica pode decorrer da construção, demonstração e avaliação de um artefato destinado a enfrentar um problema relevante em seu contexto de uso \parencite{hevnerEtAl2004designScience}. No presente estudo, o artefato de design não se restringe à interface conversacional: ele compreende o \textit{pipeline} de preparação documental, a base vetorial enriquecida por metadados legislativos, o fluxo RAG configurado no Open WebUI, o \textit{prompt} operacional, os \textit{scripts} de ingestão e avaliação e os artefatos que permitem auditoria das execuções.

O desenho do estudo também dialoga com a metodologia de \textit{Design Science Research} proposta por \textcite{peffersEtAl2007dsrm}. A identificação do problema aparece na caracterização do acesso qualificado a discursos parlamentares; os objetivos de solução são explicitados na introdução e nos critérios de sucesso; o desenho e desenvolvimento do artefato são descritos nesta seção; a demonstração e a avaliação são apresentadas nos experimentos; e a comunicação ocorre por meio deste manuscrito e dos artefatos públicos do projeto. Essa correspondência deve ser entendida como enquadramento metodológico proporcional a uma prova de conceito reprodutível, e não como validação institucional definitiva da solução.

\subsection{Arquitetura do chatbot legislativo}

A solução proposta concretiza-se em um \textit{chatbot} legislativo que utiliza a abordagem RAG para responder a consultas sobre discursos da 56\textordfeminine{} Legislatura do Senado Federal. A arquitetura compreende cinco componentes articulados: uma instância do Open WebUI, utilizada como interface de interação e orquestração do fluxo RAG; \textit{scripts} em Python empregados na extração, preparação, importação e avaliação do conteúdo documental; o ChromaDB, utilizado como banco vetorial acessado por HTTP; o Ollama, mantido como alternativa para integração com modelos locais; e a API da OpenAI, utilizada na configuração experimental validada para geração de \textit{embeddings} e respostas.

A solução foi definida para operar localmente com uso de Docker e preparada para funcionar tanto com \textit{embeddings} locais, via Ollama, quanto com \textit{embeddings} providos por serviços externos, como a API da OpenAI. Na configuração experimental validada, adotou-se esta segunda opção. Essa flexibilidade amplia a adaptabilidade da solução a diferentes restrições institucionais de infraestrutura, custo ou política tecnológica.

\subsection{Pipeline desenvolvido}

O \textit{pipeline} do projeto inicia-se com a aquisição dos dados, baseada no \textit{dataset} utilizado no projeto\footnote{\href{https://huggingface.co/datasets/fabriciosantana/discursos-senado-legislatura-56}{\textit{Dataset} público no Hugging Face}.}, que reúne discursos da 56\textordfeminine{} Legislatura do Senado Federal. Esse \textit{dataset} foi construído em trabalho anterior a partir do portal de dados abertos do Senado. Na etapa de pré-processamento e seleção textual, priorizou-se o campo ``TextoDiscursoIntegral'', recorrendo-se aos campos ``Resumo'' e ``Indexacao'' quando necessário, de modo a privilegiar o conteúdo mais completo e informativo na construção da base.

Em seguida, realizou-se o \textit{chunking}, com geração de segmentos enriquecidos por metadados relevantes\footnote{\href{https://github.com/fabriciosantana/mcdia/blob/main/05-iag/4-project/knowledge_openwebui/build_metadata.json}{Metadados públicos da base construída}.}. Esses metadados incluem data, autor, partido, unidade da federação, tipo de uso da palavra, resumo, indexação e URL do texto integral. Tal modelagem favorece não apenas a recuperação semântica, mas também a verificação posterior das respostas, uma vez que permite reconstituir a origem documental dos trechos utilizados.

A partir desses segmentos, foram gerados lotes de ingestão para o Open WebUI em formato Markdown, além de um arquivo \texttt{jsonl} de referência\footnote{\href{https://github.com/fabriciosantana/mcdia/blob/main/05-iag/4-project/knowledge_openwebui/discursos_chunks.jsonl}{Arquivo público de chunks em JSONL}.}. A importação desses lotes foi automatizada por script responsável por realizar o upload, monitorar o processamento e viabilizar a inclusão dos conteúdos na \textit{knowledge base} do Open WebUI. Concluída a ingestão, procedeu-se à indexação vetorial, com geração e persistência dos \textit{embeddings} no ChromaDB.

Na etapa de consulta, o fluxo RAG opera com recuperação vetorial, montagem do contexto recuperado, aplicação do \textit{prompt} operacional configurado no Open WebUI\footnote{\href{https://github.com/fabriciosantana/mcdia/blob/main/05-iag/4-project/eval/prompts/rag_prompt.md}{\textit{Prompt} operacional público}.} e geração da resposta pelo modelo, com parte desses passos internalizada na própria plataforma. Por fim, a avaliação foi apoiada por automação em Python, que consulta a \textit{knowledge base} via API e produz saídas nos formatos \texttt{.jsonl}, \texttt{.md} e \texttt{.csv}, viabilizando análise estruturada, avaliação por LLM como juiz, revisão qualitativa e comparação entre execuções.

\begin{figure}[ht]
\centering
\begin{tikzpicture}[
  node distance=0.72cm,
  box/.style={
    draw,
    rounded corners,
    align=center,
    minimum width=4.8cm,
    minimum height=0.82cm,
    font=\small
  },
  arrow/.style={-{Latex}, thick}
]
\node[box] (dataset) {Dataset de discursos\\56\textordfeminine{} Legislatura};
\node[box, below=of dataset] (prep) {Preparação documental\\seleção textual, chunking e metadados};
\node[box, below=of prep] (ingest) {Artefatos de ingestão\\Markdown em lotes e JSONL};
\node[box, below=of ingest] (openwebui) {Open WebUI\\ingestão e orquestração RAG};
\node[box, below=of openwebui] (vector) {Embeddings e ChromaDB\\indexação vetorial};
\node[box, below=of vector] (query) {Consulta em linguagem natural\\recuperação, contexto e geração};
\node[box, below=of query] (eval) {Avaliação\\LLM juiz, rubrica e análise manual};

\draw[arrow] (dataset) -- (prep);
\draw[arrow] (prep) -- (ingest);
\draw[arrow] (ingest) -- (openwebui);
\draw[arrow] (openwebui) -- (vector);
\draw[arrow] (vector) -- (query);
\draw[arrow] (query) -- (eval);
\end{tikzpicture}
\caption{Fluxo geral da prova de conceito RAG para consulta a discursos legislativos}
\label{fig:pipeline-rag}
\end{figure}

\subsection{Principais contribuições}

O projeto apresenta contribuições em dois níveis complementares: técnico-metodológico e institucional. No plano técnico-metodológico, o trabalho constrói uma prova de conceito reprodutível de \textit{chatbot} legislativo com RAG, adaptada a um acervo real do Senado Federal. A contribuição não consiste em propor um \textit{benchmark} definitivo para avaliação de RAG legislativo, mas em documentar um \textit{pipeline} completo, com preparação da base, enriquecimento de \textit{chunks} por metadados, ingestão, consulta, avaliação automatizada, análise qualitativa e discussão explícita dos limites observados. Nesse percurso, o estudo evidencia empiricamente a sensibilidade do desempenho à formulação das perguntas, ao controle de escopo, à escolha do LLM juiz e à configuração da recuperação semântica.

No plano institucional, a solução oferece evidência inicial de viabilidade para aplicar técnicas contemporâneas de recuperação e geração ao contexto do Parlamento brasileiro, aproximando recomendações gerais da literatura sobre RAG às condições de um acervo legislativo público já explorado por estudos computacionais sobre discursos em plenário \parencite{blanco2025plenarioPalanqueEstudio,blanco2026letrasGraudas}. Desse modo, o trabalho contribui para a discussão sobre uso responsável de inteligência artificial em contextos parlamentares e sobre formas de manter as respostas vinculadas a documentos verificáveis.

\subsection{Riscos e limitações}

Os principais riscos e limitações identificados relacionam-se à sensibilidade do desempenho à estratégia de \textit{chunking}, à possibilidade de combinação excessiva de fontes em perguntas muito abertas, à dependência da qualidade dos documentos e metadados de origem, ao risco de respostas excessivamente amplas ou inferenciais e à variabilidade da avaliação por LLM como juiz. Além disso, em cenários de restrição de recursos computacionais locais, a solução pode depender operacionalmente de serviços externos, tanto para geração de \textit{embeddings} quanto para parte da geração de respostas.

Esses riscos devem ser compreendidos em perspectiva ampliada. A literatura sobre RAG e governança pública de IA indica que sistemas apoiados em recuperação documental podem melhorar o controle sobre fontes e contexto, mas não eliminam automaticamente problemas de qualidade documental, conformidade, monitoramento ou avaliação \parencite{gao2023ragSurvey,deAlmeidaSantos2025aiGovernance,saadFalconEtAl2024ares}. Em outras palavras, o RAG reduz determinadas fragilidades do uso isolado de LLMs, mas não substitui processos institucionais de revisão.

\section{Experimentos e demonstração}

\subsection{Setup experimental}

No enquadramento metodológico adotado, os experimentos correspondem à demonstração e à avaliação inicial do artefato de design. Seu objetivo é verificar se o \textit{pipeline} proposto consegue construir uma base consultável, recuperar evidências documentais, gerar respostas fundamentadas e preservar rastreabilidade suficiente para inspeção posterior, no escopo delimitado da 56\textordfeminine{} Legislatura do Senado Federal.

O experimento utilizou uma base derivada do \textit{dataset} utilizado no projeto\footnote{\href{https://huggingface.co/datasets/fabriciosantana/discursos-senado-legislatura-56}{\textit{Dataset} público no Hugging Face}.}, contendo 15.729 registros do conjunto de origem, dos quais 15.726 correspondiam a discursos efetivamente escritos e 3 foram descartados por ausência do texto integral. A escolha desse corpus dialoga com estudos que analisam computacionalmente discursos do Senado no período 2007--2024, mas recorta a 56\textordfeminine{} Legislatura para construir e avaliar uma base RAG específica \parencite{blanco2025plenarioPalanqueEstudio,blanco2026letrasGraudas}. A partir desses registros, foram gerados 23.806 \textit{chunks}, organizados em 120 arquivos Markdown de ingestão. Esses números foram registrados nos artefatos públicos de construção da base\footnote{\href{https://github.com/fabriciosantana/mcdia/blob/main/05-iag/4-project/knowledge_openwebui/build_metadata.json}{Metadados de construção da base}; \href{https://github.com/fabriciosantana/mcdia/blob/main/05-iag/4-project/knowledge_openwebui/discursos_chunks.jsonl}{arquivo JSONL de chunks}.}, assegurando rastreabilidade quantitativa da base construída.

Na etapa de preparação da base, adotaram-se como parâmetros principais \texttt{max\_words = 850}, \texttt{overlap\_words = 150} e \texttt{chunks\_per\_file = 200}, buscando equilibrar granularidade e preservação de contexto. Além do texto segmentado, cada \textit{chunk} foi enriquecido com metadados como data, autor, partido, unidade federativa, casa legislativa, tipo de uso da palavra, resumo, indexação e URL do discurso original, de modo a favorecer recuperação qualificada e verificação das respostas. Já na etapa posterior de ingestão no Open WebUI, os arquivos \textit{Markdown} previamente segmentados foram processados com \textit{Markdown Header Text Splitter} habilitado, \textit{Chunk Size} = 6000, \textit{Chunk Overlap} = 500 e \textit{Chunk Min Size Target} = 2000, além do uso do modelo de embedding \texttt{text-embedding-3-small} e lote de \textit{embeddings} de tamanho 32.

O \textit{prompt} de RAG foi evoluído iterativamente para atender a requisitos centrais do problema institucional: privilegiar o contexto recuperado, proibir explicitamente a invenção de fatos, exigir declaração de insuficiência quando a informação solicitada não estivesse presente no contexto, distribuir citações ao longo da resposta e exigir uma seção final de referências com os metadados dos discursos utilizados. O \textit{prompt} configurado no modelo RAG do Open WebUI está disponível no repositório do projeto\footnote{\href{https://github.com/fabriciosantana/mcdia/blob/main/05-iag/4-project/eval/prompts/rag_prompt.md}{\textit{Prompt} do fluxo RAG}.} e foi utilizado tanto no uso interativo do sistema na própria plataforma quanto na bateria automatizada. Nas consultas automatizadas, o \textit{script} enviou a pergunta do usuário e anexou a coleção indexada, deixando que o Open WebUI aplicasse o \textit{template} RAG configurado no servidor e realizasse internamente a recuperação, a montagem do contexto e a formulação da resposta. Com isso, buscou-se aproximar a avaliação automatizada do fluxo de uso validado manualmente, preservando fidelidade ao contexto recuperado, controle de escopo, reconhecimento explícito de ambiguidade e estrutura de saída orientada à verificação das fontes. Tal configuração alinha-se à literatura que enfatiza verificabilidade, controle documental e avaliação explícita em sistemas RAG \parencite{gao2023ragSurvey,saadFalconEtAl2024ares,esEtAl2023ragas}.

Para a avaliação, desenvolveu-se um \textit{script} em \textit{Python} que automatizou consultas à \textit{knowledge base} do Open WebUI via API, gerando saídas em \texttt{.jsonl}, \texttt{.md} e \texttt{.csv}. O script envia cada pergunta ao endpoint de chat da plataforma e anexa a coleção indexada como \textit{file} do tipo \texttt{collection}; com isso, a recuperação e a montagem do contexto final ficam a cargo do próprio Open WebUI, e não de lógica externa implementada no \textit{script}. Os artefatos preservam a pergunta, a resposta, as fontes recuperadas, a configuração da rodada, os identificadores da \textit{knowledge base} e os \textit{hashes} dos principais arquivos usados na preparação da base. A pontuação foi atribuída em uma segunda etapa de avaliação por LLM como juiz, segundo a estratégia de \textit{LLM as a Judge}, aplicada sobre rubrica com cinco critérios (\textit{aderência}, \textit{precisão factual}, \textit{foco nas fontes}, \textit{síntese} e \textit{confiabilidade/alucinação}), cada um com escala de 0 a 2\footnote{\href{https://github.com/fabriciosantana/mcdia/blob/main/05-iag/4-project/eval/RUBRIC.md}{Rubrica de avaliação automatizada}.}. O script permite configurar separadamente o modelo usado como LLM juiz; nas rodadas principais analisadas, por simplificação experimental, optou-se por reutilizar o mesmo modelo nas etapas de geração e avaliação. Essa combinação entre LLM juiz e rubrica explícita é compatível com a literatura sobre avaliação automatizada de geração textual e de saídas de RAG \parencite{liuEtAl2025judgeRag,liuEtAl2023geval}. Nas execuções principais, o modelo usado para geração e avaliação foi \texttt{gpt-5-nano}, com recuperação sobre coleção indexada no Open WebUI, e os \textit{embeddings} \texttt{text-embedding-3-small} permaneceram registrados nos metadados dos trechos retornados.

Assim, a avaliação combina quatro camadas complementares, mas não equivalentes. A recuperação é observada por meio dos metadados dos trechos e arquivos retornados pelo Open WebUI; a geração é examinada nas respostas completas produzidas pelo fluxo RAG; a avaliação automatizada atribui notas às respostas por meio de LLM como juiz e rubrica explícita; e a inspeção humana/manual atua como verificação qualitativa amostral de casos fortes, cautelosos, limítrofes e divergentes. Essa separação é importante porque uma nota alta na resposta final não demonstra, isoladamente, que a recuperação foi ótima, nem que o julgamento automatizado substitui validação humana independente.

Do ponto de vista da recuperação, o \textit{script} de avaliação não fixa explicitamente um \textit{top-k} nem aplica reranqueamento manual: esses passos permanecem internalizados na plataforma. Ainda assim, os artefatos da rodada base permitem observar o comportamento efetivo do sistema\footnote{\href{https://github.com/fabriciosantana/mcdia/blob/main/05-iag/4-project/eval/results/rag_eval_20260416T172816Z.question_analysis.md}{Análise da rodada base por pergunta}.}. Nessa rodada, cada resposta retornou metadados associados a 9 a 15 trechos-fonte, com média de 12,1, distribuídos por 3 a 13 arquivos \textit{Markdown} de origem, com média de 8,6 arquivos por resposta. Isso sugere que o contexto final utilizado na geração resultou da agregação automática de múltiplos trechos recuperados semanticamente a partir da coleção indexada.

A tipologia das perguntas não foi tomada integralmente de um único \textit{benchmark}, mas foi inspirada por \textit{benchmarks} de avaliação de RAG, especialmente o RAGEval, que organiza consultas em famílias como pergunta factual, integração de informações, comparação numérica, sequência temporal, raciocínio em múltiplas etapas, sumarização e perguntas não respondíveis \parencite{zhuEtAl2025rageval}. A partir dessa base, o presente experimento adaptou as categorias ao domínio parlamentar e aos riscos específicos observados no \textit{corpus} de discursos do Senado. Assim, categorias como \textit{factual}, \textit{comparação}, \textit{síntese}, \textit{multietapas} e \textit{evidência negativa} dialogam diretamente com a literatura, enquanto rótulos como \textit{foco no autor}, \textit{desambiguação}, \textit{autoria cruzada}, \textit{controle de escopo} e \textit{estresse de recuperação} devem ser entendidos como extensões metodológicas propostas neste trabalho, criadas para tensionar problemas especialmente relevantes em consultas sobre autoria, mistura indevida de discursos, controle de escopo e robustez da recuperação.

\subsection{Resultados alcançados}

Os resultados podem ser descritos em três frentes complementares: integridade da base, inspeção manual de respostas e bateria automatizada de perguntas. 

Na frente operacional, a base foi carregada e manteve-se operacional nas consultas realizadas, em coerência com os metadados gerados na preparação: 15.729 registros de entrada, 15.726 discursos aproveitados, 3 descartes por ausência de texto útil e 23.806 \textit{chunks} persistidos para recuperação semântica. Esses números não demonstram qualidade da resposta por si só, mas indicam que o experimento partiu de uma base operacionalmente íntegra e rastreável.

Na inspeção manual, adotaram-se três estudos de caso exemplificativos, documentados no relatório público de resultados manuais\footnote{\href{https://github.com/fabriciosantana/mcdia/blob/main/05-iag/4-project/eval-openwebui/RESULTS.md}{Relatório público de resultados manuais}.}, escolhidos para representar um caso forte, um caso de cautela e um caso de limite. No primeiro, a resposta sobre os argumentos de Paulo Paim contra a capitalização da Previdência mostrou desempenho verificável consistente: o sistema recuperou múltiplos discursos do senador, sintetizou os principais pontos e apresentou referências explícitas ao final. No segundo, diante de uma pergunta de desambiguação de autoria sobre Confúcio Moura, o sistema adotou postura cautelosa, reconhecendo insuficiência de contexto e evitando atribuições indevidas. No terceiro, diante de uma pergunta multifocal que exigia reconstrução de um plano legislativo consolidado para reduzir desigualdade, o sistema preservou parte dessa cautela, mas já operou em nível de abstração mais alto, revelando o limite do fluxo quando a consulta exige consolidação ampla, articulação entre autores e fechamento programático a partir de evidências fragmentárias.

A avaliação automatizada foi consolidada a partir das três rodadas completas executadas após o alinhamento entre o fluxo automatizado e a configuração operacional do Open WebUI\footnote{\href{https://github.com/fabriciosantana/mcdia/blob/main/05-iag/4-project/eval/results/rag_eval_20260416T172816Z.csv}{CSV da rodada base}; \href{https://github.com/fabriciosantana/mcdia/blob/main/05-iag/4-project/eval/results/rag_eval_20260416T182240Z.csv}{CSV da segunda rodada}; \href{https://github.com/fabriciosantana/mcdia/blob/main/05-iag/4-project/eval/results/rag_eval_20260416T232921Z.csv}{CSV da terceira rodada}; \href{https://github.com/fabriciosantana/mcdia/blob/main/05-iag/4-project/eval/results/stability_summary_20260417T012001Z.md}{síntese de estabilidade}.}. Em cada rodada, 20 perguntas foram submetidas à \textit{knowledge base}, todas as respostas foram geradas com sucesso e, em seguida, avaliadas em uma segunda etapa de avaliação por LLM como juiz, com base na rubrica definida para o experimento. As médias observadas foram 9,15, 9,55 e 9,35 pontos em 10, resultando em média agregada de 9,35. Não houve falhas nas três execuções. No vocabulário da rubrica adotada, 17 respostas da primeira rodada ficaram na faixa ``pronta para uso'' (8--10) e 3 na faixa ``utilizável com supervisão'' (5--7); nas duas rodadas seguintes, 19 respostas ficaram na primeira faixa e 1 na segunda. Essas faixas indicam qualidade estimada no escopo do experimento, e não prontidão operacional ampla. O conjunto de 20 perguntas, suas categorias e as notas obtidas nas três rodadas está disponível nos artefatos públicos da avaliação\footnote{\href{https://github.com/fabriciosantana/mcdia/blob/main/05-iag/4-project/eval/discursos_questions.json}{Conjunto de perguntas}; \href{https://github.com/fabriciosantana/mcdia/blob/main/05-iag/4-project/eval/results/rag_eval_20260416T172816Z.question_analysis.md}{análise por pergunta da rodada base}.}. Esse material já distingue categorias inspiradas diretamente na literatura de \textit{benchmarks} RAG \parencite{zhuEtAl2025rageval} e categorias introduzidas como adaptação metodológica deste trabalho.

\begin{table}[t]
\centering
\footnotesize
\caption{Síntese entre objetivos, evidências e interpretação dos resultados}
\label{tab:sintese-objetivos-resultados}
\begin{tabular}{@{}>{\raggedright\arraybackslash}p{0.24\textwidth}>{\raggedright\arraybackslash}p{0.38\textwidth}>{\raggedright\arraybackslash}p{0.30\textwidth}@{}}
\textbf{Objetivo avaliado} & \textbf{Evidência empírica} & \textbf{Interpretação} \\
\hline
Construir base verificável & 15.726 discursos aproveitados, 23.806 chunks gerados e metadados legislativos preservados & A base construída mostrou-se operacional e rastreável para a avaliação realizada. \\
Validar consulta assistida & Três estudos de caso manuais representando caso forte, caso de cautela e caso de limite & O sistema apresenta melhor desempenho quando a pergunta é delimitada por autor, tema ou evidência documental. \\
Avaliar desempenho automatizado & Três rodadas completas, com médias 9,15, 9,55 e 9,35, sem falhas operacionais & O desempenho agregado foi alto nas métricas adotadas, embora com variação em itens mais complexos. \\
Identificar limites de uso & Oscilação em perguntas de precisão numérica, controle de escopo, multietapas e autoria cruzada & Consultas de maior abstração exigem supervisão humana e formulação cuidadosa. \\
Verificar robustez da avaliação por LLM & Análise manual amostral com convergência geral, mas divergência relevante em item de checagem de alucinação & O LLM como juiz é útil para escalar a avaliação, mas não substitui validação humana em casos críticos. \\
\hline
\end{tabular}
\end{table}

Rodadas complementares com outro modelo na etapa de avaliação por LLM foram usadas apenas como teste. Embora tenham completado a bateria, elas atribuíram nota máxima a todos os itens, o que sugeriu um padrão excessivamente permissivo de julgamento por parte do LLM juiz alternativo. Por essa razão, tais execuções não foram adotadas como evidência quantitativa principal, mas como alerta metodológico sobre a influência do modelo avaliador na estratégia de \textit{LLM as a Judge}.

Os resultados por família de perguntas e por tipo de exigência cognitiva ajudam a qualificar melhor o comportamento do sistema, embora esta bateria ainda seja pequena para caracterizar de modo estável todas as categorias consideradas. Nesta amostra de 20 perguntas, consultas das categorias \textit{factual}, \textit{proposta}, \textit{foco no autor}, \textit{desambiguação} e \textit{evidência negativa} mantiveram desempenho alto e pouca variação entre as rodadas, o que sugere boa aderência ao contexto quando a tarefa exigia localizar autoria, sintetizar argumentos delimitados ou reconhecer ausência de evidência suficiente. Em contrapartida, as perguntas que mais tensionaram o sistema foram aquelas que exigiam precisão numérica, controle de escopo, reconstrução multietapas, comparação entre autores e distinção fina entre formulações semelhantes. A maior variação ocorreu na pergunta \texttt{q15}\footnote{\href{https://github.com/fabriciosantana/mcdia/blob/main/05-iag/4-project/eval/discursos_questions.json}{Conjunto público de perguntas da bateria automatizada}.}, voltada à recuperação de informação numérica específica, cujas notas foram 6, 10 e 10; também apresentaram oscilação relevante a \texttt{q16}, que tensionava controle de escopo, com notas 7, 9 e 10, e a \texttt{q18}, formulada para testar autoria cruzada, com notas 8, 7 e 9. Esses casos sugerem que parte da instabilidade não decorre de falha operacional, mas da dificuldade intrínseca de avaliar respostas que combinam múltiplas fontes, números dispersos, inferências de autoria e limites de extrapolação.

Esses achados convergem com a avaliação qualitativa do projeto: perguntas mais delimitadas, ancoradas em um autor, tema ou argumento específico, tendem a produzir respostas mais fortes e verificáveis; perguntas amplas, multifocais ou formuladas de modo a induzir generalização aumentam o risco de mistura indevida de discursos, atribuições excessivas e perda de precisão factual. A análise manual amostral da rodada base, conduzida sobre sete respostas selecionadas, indicou convergência com o LLM juiz na maior parte dos casos, mas também revelou uma divergência importante em pergunta de checagem de alucinação, na qual a avaliação humana considerou a resposta inadequada por erro de escopo apesar da nota máxima atribuída automaticamente. Essa divergência não deve ser tratada apenas como ruído avaliativo: ela constitui achado metodológico relevante, pois mostra que o juiz automatizado pode superestimar respostas que parecem formalmente adequadas, mas violam o escopo documental esperado. Essa triangulação reforça que a avaliação por LLM como juiz é útil para escalar a análise, mas deve ser interpretada como instrumento auxiliar, especialmente em tarefas abertas ou metavaliativas. Em termos metodológicos, isso sugere que a utilidade da prova de conceito é mais clara para consulta assistida e exploração do acervo com apoio documental, mas que respostas de maior abstração ainda dependem de supervisão humana e de refinamentos adicionais no desenho de recuperação, no controle de escopo e na formulação do prompt \parencite{gao2023ragSurvey,wangEtAl2025segmentation,saadFalconEtAl2024ares}.

A inclusão de referências documentais ao longo da resposta e ao seu final mostrou-se viável por meio de instruções adequadas no \textit{prompt}, o que facilita a conferência posterior das informações. Em síntese, os resultados sustentam a viabilidade técnica inicial da prova de conceito e oferecem evidência mais concreta de onde o sistema apresenta utilidade e onde seus limites permanecem mais visíveis. Para tornar a demonstração menos abstrata, os exemplos completos de interação manual estão disponíveis em relatório público do projeto\footnote{\href{https://github.com/fabriciosantana/mcdia/blob/main/05-iag/4-project/eval-openwebui/RESULTS.md}{Relatório público de resultados manuais}.}, incluindo um caso de bom desempenho, um caso em que o sistema reconhece insuficiência de contexto e evita atribuições indevidas de autoria, e um caso mais desafiador em que a pergunta exige consolidação legislativa ampla sob forte restrição de contexto.

\subsection{Ameaças à validade}

Os resultados devem ser interpretados como evidência inicial de uma prova de conceito reprodutível, e não como \textit{benchmark} definitivo para sistemas RAG legislativos. A principal ameaça à validade interna decorre do tamanho reduzido da bateria automatizada, composta por 20 perguntas, e do balanceamento desigual entre categorias: algumas famílias aparecem com poucos exemplos, o que limita inferências estáveis por tipo de pergunta. Há também ameaça associada ao uso de \textit{LLM as a Judge}, pois os testes de sensibilidade indicaram que a pontuação pode variar de forma relevante conforme o modelo avaliador e a rubrica utilizada. A validade de construto é limitada pelo fato de a avaliação combinar notas automatizadas, inspeção manual e estudos de caso, sem ainda dispor de validação humana ampla e independente. Além disso, esta etapa não incluiu baseline lexical ou sem RAG, estudo de ablação dos metadados, métricas formais de recuperação nem protocolo independente de anotação humana; tais procedimentos ficam como prioridades para uma rodada experimental posterior. Por fim, a validade externa deve ser considerada com cautela, pois os resultados dependem da configuração específica do Open WebUI, do \textit{prompt} operacional, dos parâmetros de \textit{chunking}, da coleção indexada e do corpus de discursos da 56\textordfeminine{} Legislatura. Assim, os achados são mais bem compreendidos como demonstração reprodutível de viabilidade e limites em um domínio parlamentar específico, passível de replicação e expansão em estudos posteriores.

\subsection{Potencial de generalização}

Embora o experimento tenha sido desenhado para discursos da 56\textordfeminine{} Legislatura do Senado Federal, parte do pipeline tende a ser reaproveitável em outros acervos legislativos. São potencialmente transferíveis a preparação documental em lotes, o enriquecimento dos \textit{chunks} com metadados, a ingestão em base vetorial, o uso de \textit{prompt} orientado a evidências, o registro de artefatos de execução e a avaliação por LLM como juiz combinada à análise manual amostral. Por outro lado, permanecem específicos do \textit{corpus} analisado a estrutura dos campos de origem, os metadados disponíveis, a natureza dos pronunciamentos em plenário, a distribuição temporal dos discursos e a tipologia de perguntas construída para tensionar problemas de autoria, escopo e recuperação nesse domínio. Assim, a contribuição mais generalizável do trabalho não é um conjunto universal de resultados, mas um protocolo reprodutível e adaptável para construção e avaliação inicial de sistemas RAG em acervos legislativos.

A transferência do protocolo para outros acervos exigiria, portanto, nova caracterização documental e institucional. Diferenças de idioma, formato dos documentos, estabilidade das URLs, granularidade dos metadados, regime de dados abertos, atualização do corpus e infraestrutura disponível podem alterar tanto a recuperação quanto a verificabilidade das respostas. Em bases legislativas com retificações frequentes, documentos substituídos, discursos complementados posteriormente ou campos incompletos, a qualidade do sistema dependeria de regras explícitas de versionamento, descarte, atualização e auditoria da base.

\subsection{Impacto na administração pública}

O impacto potencial do \textit{chatbot} legislativo com RAG na administração pública pode ser analisado em pelo menos quatro dimensões. Em primeiro lugar, há potencial de contribuição para a transparência, na medida em que as respostas passam a ser ancoradas em discursos públicos e citáveis, facilitando a verificação por terceiros. Em segundo lugar, observa-se potencial de ganho de eficiência, uma vez que o sistema pode reduzir o tempo necessário para localizar informações relevantes em grandes acervos documentais. Em terceiro lugar, a solução pode contribuir para a memória institucional, ao facilitar a recuperação de posicionamentos anteriores, temas recorrentes e evidências documentais. Por fim, identifica-se potencial de apoio à decisão, na medida em que gabinetes e equipes técnicas podem dispor de uma camada adicional de mediação informacional para explorar o acervo legislativo de forma mais estruturada \parencite{martelloteViola2025organizacaoConhecimento,deAlmeidaSantos2025aiGovernance,wfd2023guidelinesAiParliaments}.

Por operar sobre fontes controladas e permitir referência explícita aos discursos da 56\textordfeminine{} Legislatura, a solução é mais compatível com exigências de \textit{accountability} do que o uso isolado de um \textit{chatbot} generalista. Essa característica assume especial relevância em contextos públicos, nos quais a confiabilidade da informação não pode ser dissociada da governança das fontes e da possibilidade de auditoria posterior.

Ainda assim, o potencial institucional identificado não equivale a prontidão operacional. Em eventual adoção por órgão público, seria necessário definir responsabilidades de curadoria da base, periodicidade de atualização, tratamento de discursos removidos ou retificados, trilhas de auditoria, política de retenção de logs, controle de versões do \textit{prompt} e mecanismos de contestação ou correção de respostas. Também seria necessário distinguir usos assistivos, voltados à exploração documental por servidores e pesquisadores, de usos voltados ao atendimento público, que exigiriam níveis mais altos de validação, supervisão humana e governança.

\section{Conclusão e trabalhos futuros}

\subsection{Síntese das contribuições}

Neste trabalho, foi desenvolvida uma prova de conceito reprodutível de \textit{chatbot} legislativo com RAG voltada ao acervo de discursos da 56\textordfeminine{} Legislatura do Senado Federal. A solução compreende um \textit{pipeline} de preparação documental, uma base vetorial, uma interface de consulta e \textit{scripts} auxiliares de ingestão, avaliação e registro dos resultados. Demonstrou-se, assim, que é possível estruturar uma base semântica consultável e responder perguntas em linguagem natural com apoio em evidências recuperadas, aproximando o uso de LLMs das exigências de atualidade, verificabilidade e controle das fontes apontadas pela literatura \parencite{gao2023ragSurvey}. A contribuição é, portanto, metodológica e aplicada, pois não se limita ao artefato implementado: envolve também um protocolo reprodutível, preservação de metadados legislativos e uma estratégia explícita de avaliação dos resultados.

Sob o enquadramento de \textit{Design Science Research} adotado no artigo, essa contribuição deve ser entendida como construção e avaliação inicial de um artefato de design para mediação informacional em acervo parlamentar, e não como implantação operacional validada em ambiente institucional.

\subsection{Discussão dos resultados frente aos objetivos}

À luz dos objetivos propostos, os resultados constituem evidência inicial de viabilidade. A prova de conceito atendeu aos critérios centrais definidos para indexação da base de discursos da 56\textordfeminine{} Legislatura, recuperação de trechos relevantes, geração de respostas fundamentadas e avaliação estruturada. As três rodadas automatizadas completaram a bateria de 20 perguntas sem falhas operacionais e apresentaram médias entre 9,15 e 9,55 em escala de 0 a 10. A inspeção manual e a triangulação amostral reforçaram, contudo, que a qualidade de um sistema RAG institucional depende de disciplina metodológica e de decisões arquiteturais consistentes. Estratégias de \textit{chunking}, desenho do \textit{prompt}, qualidade dos metadados, escolha do LLM juiz e avaliação iterativa influenciam diretamente o comportamento do \textit{chatbot}.

Tal constatação é compatível com a literatura, segundo a qual RAG não deve ser entendido apenas como a conexão entre um LLM e um banco vetorial, mas como uma solução sociotécnica que exige avaliação contínua, rastreabilidade, governança e critérios explícitos de sucesso, especialmente quando utilizada em funções públicas sensíveis \parencite{gao2023ragSurvey,deAlmeidaSantos2025aiGovernance,saadFalconEtAl2024ares}. Nesse sentido, a principal contribuição do trabalho não reside apenas na implementação do protótipo, mas também na explicitação das escolhas técnicas, limitações práticas e condições institucionais envolvidas em sua construção.

\subsection{Limitações}

Entre as limitações atuais da solução, destacam-se o tamanho ainda reduzido da bateria automatizada, a necessidade de validação humana mais ampla, a ausência, nesta etapa, de filtros mais sofisticados por autor, período ou tipo de discurso diretamente integrados ao fluxo de geração, a possibilidade de respostas excessivamente amplas em consultas abertas e a dependência de serviços externos para embeddings e parte da geração. Além disso, testes com LLM juiz alternativo indicaram que a estratégia de \textit{LLM as a Judge} é sensível ao modelo escolhido, podendo produzir avaliações excessivamente permissivas. Tais limitações dialogam com alertas da literatura acerca de riscos operacionais, vieses de dados, segmentação documental e necessidade de monitoramento contínuo em sistemas RAG \parencite{gao2023ragSurvey,wangEtAl2025segmentation,liuEtAl2025judgeRag}.

Há também limitações institucionais que não se resolvem apenas por ajuste técnico do pipeline. A qualidade de um sistema desse tipo depende de curadoria contínua do acervo, atualização controlada da base, definição de responsáveis por validação documental e regras para lidar com registros corrigidos, removidos ou incompletos. Por isso, os resultados devem ser lidos como evidência de viabilidade para consulta assistida e pesquisa aplicada, não como autorização para uso autônomo em processos decisórios ou atendimento público sem governança adicional.

\subsection{Trabalhos futuros}

Como trabalhos futuros, recomenda-se ampliar a bateria de avaliação, balancear melhor as categorias de perguntas, comparar avaliações por LLM como juiz com validações humanas independentes e experimentar \textit{frameworks} especializados, como RAGAS, para complementar a rubrica utilizada na avaliação por LLM. Recomenda-se, ainda, incorporar filtros estruturados por metadado na etapa de recuperação, experimentar reranqueadores mais robustos, testar diferentes estratégias de \textit{chunking} por tipo documental e implementar pós-processamento determinístico para a seção de referências.

Além disso, a prova de conceito pode ser adaptada para outros tipos de documentos legislativos, como pareceres, projetos de lei, notas técnicas e documentos administrativos relacionados ao processo legislativo. Essa expansão exigiria reavaliar metadados, estratégias de \textit{chunking}, categorias de pergunta, rubrica e critérios de validação em função de cada novo acervo. Do ponto de vista institucional, trabalhos futuros também devem detalhar protocolos de atualização da base, papéis de curadoria, validação humana, monitoramento de erros, registro de interações e critérios para separar respostas exploratórias de respostas aptas a circular em fluxos formais. Em horizonte mais avançado, um desdobramento possível consiste na evolução do sistema para cenários de geração assistida de minutas, relatórios ou discursos, desde que preservados mecanismos de verificação documental, supervisão humana obrigatória e controles institucionais adequados. Dessa forma, o presente trabalho abre caminho para agenda mais ampla de pesquisa e desenvolvimento sobre inteligência artificial generativa responsável no contexto do Parlamento brasileiro.

\section{Declarações de integridade científica e ciência aberta}

\subsection{Conflitos de interesse e financiamento}

Os autores declaram não possuir conflitos de interesse relacionados ao desenvolvimento, à avaliação ou à apresentação dos resultados deste trabalho. A prova de conceito foi desenvolvida em contexto acadêmico, sem financiamento externo específico declarado para esta pesquisa.

\subsection{Disponibilidade de dados, código e artefatos}

Os dados utilizados derivam de registros públicos de discursos parlamentares do Senado Federal e foram organizados em \textit{dataset} público no Hugging Face\footnote{\href{https://huggingface.co/datasets/fabriciosantana/discursos-senado-legislatura-56}{\textit{Dataset} público de discursos da 56\textordfeminine{} Legislatura}.}. O código, os \textit{scripts} auxiliares, os parâmetros de preparação da base, a rubrica, o \textit{prompt} operacional e os artefatos experimentais citados no artigo estão disponíveis no repositório público do projeto\footnote{\href{https://github.com/fabriciosantana/mcdia/tree/main/05-iag/4-project}{Repositório público do projeto}.}. Essa disponibilização busca favorecer auditoria, reprodutibilidade e reutilização crítica do protocolo em estudos posteriores.

\subsection{Uso de inteligência artificial generativa}

Ferramentas de inteligência artificial generativa foram utilizadas como apoio à programação, à revisão textual, à organização de argumentos e à avaliação automatizada descrita na metodologia. As decisões de escopo, interpretação dos resultados, seleção das evidências, redação final e responsabilidade pelo conteúdo permanecem sob responsabilidade dos autores. As saídas geradas por modelos foram revisadas criticamente e não substituem validação humana independente.

\subsection{Considerações éticas}

O estudo utiliza documentos parlamentares públicos, sem coleta direta de dados junto a participantes humanos e sem acesso a informações sigilosas. Ainda assim, por tratar de acervo institucional e de respostas geradas por IA, o trabalho adota cautela na interpretação dos resultados e enfatiza a necessidade de verificação documental, supervisão humana e governança adequada antes de qualquer uso operacional em contexto público.

\input{sections/08-referencias}
% Apêndices mantidos no repositório, mas omitidos da versão atual do artigo.
% Os artefatos correspondentes são referenciados por notas de rodapé no texto.
% \appendix
\clearpage
\phantomsection
\refstepcounter{section}
\addcontentsline{toc}{section}{APÊNDICE \Alph{section} -- SÍNTESE DA AVALIAÇÃO AUTOMATIZADA}
\begin{center}
\textbf{APÊNDICE \Alph{section} -- SÍNTESE DA AVALIAÇÃO AUTOMATIZADA}
\end{center}
\label{sec:apendice-avaliacao}

Este apêndice reúne a síntese das três rodadas automatizadas utilizadas para complementar a demonstração empírica apresentada na seção de experimentos, com o objetivo de tornar visíveis, em formato legível, os resultados quantitativos após o alinhamento entre a automação e a configuração operacional do Open WebUI. As categorias apresentadas na tabela combinam tipos de consulta inspirados em benchmarks de avaliação de RAG, especialmente o RAGEval \parencite{zhuEtAl2025rageval}, com extensões metodológicas definidas no próprio experimento para capturar riscos típicos do domínio parlamentar.

\subsection{Síntese das rodadas automatizadas}

A Tabela~\ref{tab:rodadas-avaliacao} apresenta a síntese das três rodadas principais consideradas na análise. Em todas elas, 20 perguntas foram avaliadas com estratégia de \textit{LLM as a Judge} \parencite{liuEtAl2025judgeRag}, mantendo a mesma base de conhecimento, o mesmo conjunto de perguntas e o mesmo modelo nas etapas de geração e avaliação. A nota total de cada item corresponde à soma dos cinco critérios da rubrica, em escala de 0 a 10 \parencite{liuEtAl2023geval}.

\begingroup
\footnotesize
\begin{longtable}{@{}>{\raggedright\arraybackslash}p{0.26\textwidth}cccccc@{}}
\caption{Síntese das três rodadas principais de avaliação automatizada}
\label{tab:rodadas-avaliacao}\\
\textbf{Rodada} & \textbf{Perguntas} & \textbf{Falhas} & \textbf{Média} & \textbf{Mín.} & \textbf{8--10} & \textbf{5--7} \\
\hline
\endfirsthead
\textbf{Rodada} & \textbf{Perguntas} & \textbf{Falhas} & \textbf{Média} & \textbf{Mín.} & \textbf{8--10} & \textbf{5--7} \\
\hline
\endhead
Primeira & 20 & 0 & 9,15 & 6 & 17 & 3 \\
Segunda & 20 & 0 & 9,55 & 7 & 19 & 1 \\
Terceira & 20 & 0 & 9,35 & 7 & 19 & 1 \\
\hline
\end{longtable}
\endgroup

A Tabela~\ref{tab:resultado-final-20} detalha as notas obtidas por pergunta nas três rodadas. Os resultados completos estão disponíveis nos arquivos públicos da primeira\footnote{\url{https://github.com/fabriciosantana/mcdia/blob/main/05-iag/4-project/eval/results/rag_eval_20260416T172816Z.csv}}, segunda\footnote{\url{https://github.com/fabriciosantana/mcdia/blob/main/05-iag/4-project/eval/results/rag_eval_20260416T182240Z.csv}} e terceira\footnote{\url{https://github.com/fabriciosantana/mcdia/blob/main/05-iag/4-project/eval/results/rag_eval_20260416T232921Z.csv}} rodadas.

\begingroup
\footnotesize
\setlength{\LTpre}{0pt}
\setlength{\LTpost}{0pt}
\begin{longtable}{@{}p{0.07\textwidth}>{\raggedright\arraybackslash}p{0.15\textwidth}>{\raggedright\arraybackslash}p{0.52\textwidth}ccc@{}}
\caption{Notas por pergunta nas três rodadas automatizadas principais}\label{tab:resultado-final-20}\\
\textbf{ID} & \textbf{Categoria} & \textbf{Pergunta} & \textbf{R1} & \textbf{R2} & \textbf{R3} \\
\hline
\endfirsthead
\textbf{ID} & \textbf{Categoria} & \textbf{Pergunta} & \textbf{R1} & \textbf{R2} & \textbf{R3} \\
\hline
\endhead
q01 & factual & Quais dados Ciro Nogueira apresenta para criticar a concentração bancária no Brasil e qual conclusão ele tira sobre juros, tarifas e spread? & 10 & 10 & 9 \\
q02 & factual & Quais são os principais argumentos de Paulo Paim contra a capitalização da Previdência e que exemplos ele usa para sustentar essa posição? & 10 & 9 & 9 \\
q03 & proposta & Que propostas Confúcio Moura apresenta para melhorar a educação? & 10 & 10 & 10 \\
q04 & proposta & Por que Confúcio Moura defende que senadores e deputados acompanhem mais de perto escolas públicas? & 10 & 10 & 10 \\
q05 & foco no autor & Quais prioridades Chico Rodrigues afirma defender no início do mandato e como ele relaciona ordem, união e desenvolvimento para o Brasil e para Roraima? & 10 & 10 & 10 \\
q06 & comparação & Compare as preocupações de Paulo Paim sobre Previdência com as preocupações de Confúcio Moura sobre educação. & 10 & 10 & 10 \\
q07 & síntese & Quais temas sociais e econômicos aparecem com mais frequência nos discursos carregados e quais senadores associam esses problemas a propostas concretas de solução? & 9 & 10 & 9 \\
q08 & números & Quais discursos citam percentuais, valores ou estatísticas e para sustentar quais argumentos? & 8 & 10 & 10 \\
q09 & foco no autor & Que argumentos Wellington Fagundes usa para defender investimentos em infraestrutura e qual relação ele faz entre logística e competitividade? & 10 & 10 & 8 \\
q10 & ampla & Quais diagnósticos sobre desigualdade, vulnerabilidade ou falta de desenvolvimento aparecem nesses discursos e que tipos de solução são defendidos? & 10 & 9 & 9 \\
q11 & comparação & Compare como Ciro Nogueira e Wellington Fagundes usam dados econômicos ou setoriais para defender mudanças estruturais. Quais números aparecem, para sustentar que tese, e onde os argumentos deles divergem? & 8 & 9 & 9 \\
q12 & desambiguação & Quais propostas sobre educação aparecem especificamente em falas de Confúcio Moura e quais aparecem em falas de outros senadores, mas poderiam ser confundidas com as dele? Responda separando com clareza o que é dele e o que não é. & 10 & 10 & 10 \\
q13 & evidência negativa & Há, no contexto recuperado, evidência suficiente para afirmar que existe uma proposta detalhada de implementação, com prazos e fonte de financiamento, para combater a desigualdade? Se não houver, explique exatamente o que falta. & 10 & 10 & 10 \\
q14 & multietapas & A partir de diferentes discursos, reconstrua a cadeia de raciocínio que liga desigualdade, desenvolvimento regional, educação e proteção social. Quais autores fazem essas conexões de modo mais explícito e com quais limites? & 7 & 8 & 7 \\
q15 & precisão numérica & Quais números citados nos discursos parecem centrais para o argumento de cada autor e quais aparecem apenas como apoio periférico? Distingua os dois grupos e evite misturar estatísticas de autores diferentes. & 6 & 10 & 10 \\
q16 & controle de escopo & O que exatamente Paulo Paim afirma sobre capitalização da Previdência, e o que seria uma extrapolação indevida para além do que os trechos permitem concluir? & 7 & 9 & 10 \\
q17 & foco na fonte & Quais respostas exigiriam mais cautela por dependerem de inferências amplas a partir de poucos trechos: as sobre educação, as sobre desigualdade ou as sobre infraestrutura? Justifique com base nas limitações do próprio contexto. & 10 & 10 & 10 \\
q18 & autoria cruzada & Quais semelhanças de linguagem ou tema entre Paulo Paim, Confúcio Moura e outros autores podem levar o sistema a confundir autoria ou proposta? Identifique os riscos concretos de confusão e ilustre com exemplos. & 8 & 7 & 9 \\
q19 & estresse de recuperação & Se a pergunta for ``quais senadores apresentam diagnóstico e solução no mesmo discurso, sem depender de síntese entre vários discursos?'', quem realmente se encaixa nisso no contexto disponível? Cite apenas os casos sustentados com segurança. & 10 & 10 & 8 \\
q20 & checagem de alucinação & Entre as perguntas deste conjunto, quais tenderiam a induzir alucinação ou excesso de confiança se o modelo não reconhecesse limites do contexto? Explique por que e indique como uma resposta cautelosa deveria se comportar. & 10 & 10 & 10 \\
\hline
\end{longtable}
\endgroup

\subsection{Rubrica utilizada na avaliação automatizada}

Para favorecer reprodutibilidade, a seguir registra-se o conteúdo da rubrica aplicada pelo LLM juiz. Sua formulação em critérios explícitos e escala discreta aproxima o procedimento adotado de abordagens de avaliação por formulário e rubrica descritas na literatura de \textit{LLM-based evaluation} \parencite{liuEtAl2023geval}. As faixas de decisão registradas na rubrica têm sentido operacional no escopo do experimento, não devendo ser lidas como certificação de prontidão para uso institucional amplo.

\begin{Verbatim}[breaklines=true]
# RAG Evaluation Rubric

Use esta rubric para classificar cada resposta em uma planilha ou revisao manual.

## Escala

- `2`: bom
- `1`: aceitavel
- `0`: ruim

## Criterios

### 1. Aderencia a pergunta
- `2`: responde exatamente ao que foi perguntado
- `1`: responde parcialmente ou abre demais o escopo
- `0`: foge do foco

### 2. Precisao factual
- `2`: fatos, numeros, autores e temas corretos
- `1`: pequenos desvios ou formulacoes muito amplas
- `0`: erros factuais relevantes

### 3. Foco nas fontes
- `2`: usa trechos coerentes e pertinentes
- `1`: usa fontes validas, mas mistura demais ou inclui evidencias perifericas
- `0`: cita fontes incompatíveis ou aparentemente sem relacao

### 4. Sintese
- `2`: bem organizada, clara e proporcional
- `1`: correta, mas excessivamente longa, repetitiva ou burocratica
- `0`: confusa ou mal estruturada

### 5. Confianca / Alucinacao
- `2`: nao inventa e sinaliza limites quando necessario
- `1`: faz inferencias amplas, mas sem erros graves
- `0`: aparenta inventar ou extrapolar sem base

## Decisao pratica

- `8 a 10`: pronta para uso
- `5 a 7`: utilizavel com supervisao
- `0 a 4`: precisa ajuste de prompt, retrieval ou pergunta

## Padroes comuns de erro

- Mistura excessiva de autores ou discursos
- Resposta mais ampla que a pergunta
- Generalizacao sem apoio suficiente
- Falta de citacoes ou citacoes mal distribuídas
- Nao reconhecer quando faltou informacao no contexto
\end{Verbatim}

% \clearpage
\phantomsection
\refstepcounter{section}
\addcontentsline{toc}{section}{APÊNDICE \Alph{section} -- EXEMPLOS DE INTERAÇÃO MANUAL E DISCUSSÃO QUALITATIVA}
\begin{center}
\textbf{APÊNDICE \Alph{section} -- EXEMPLOS DE INTERAÇÃO MANUAL E DISCUSSÃO QUALITATIVA}
\end{center}
\label{sec:apendice-interacao}

Este apêndice reúne exemplos completos de interação manual utilizados para complementar a demonstração qualitativa apresentada na seção de experimentos. O objetivo é tornar auditáveis, em formato legível, um caso de bom desempenho em consulta temática delimitada, um caso em que o próprio sistema reconhece limites do contexto e evita atribuições indevidas de autoria, e um caso mais desafiador, no qual a pergunta exige consolidação legislativa ampla, articulação entre autores, cronograma, orçamento e consenso.

\subsection{Caso 1: resposta auditável em consulta temática delimitada}\label{sec:anexo-caso-bom}

\textbf{Pergunta}

Quais são os principais argumentos de Paulo Paim contra a capitalização da Previdência e que exemplos ele usa para sustentar essa posição?

\textbf{Resposta}

\textbf{Principais argumentos de Paulo Paim contra a capitalização da Previdência}

\begin{enumerate}
  \item \textbf{``Capitalização'' equivale, na prática, a ``poupança individual'', o que fragiliza a proteção previdenciária.}

  Ele afirma que ``regime de capitalização me preocupa muito'' porque ``significa uma poupança individual'', em que ``quem vai cuidar do meu dinheiro é aquele fundo'' e ``se o fundo aplicar mal, o resultado é zero'' em rendimento.

  \item \textbf{Perde-se a lógica solidária, e sobra apenas uma parcela do trabalhador.}

  Ele argumenta que se retiram diversas fontes de financiamento e se deixa ``só 10\% do empregado'', dizendo que ``claro que não vai dar'', pois ``tu tens 100, tira 90''.

  \item \textbf{Com esse modelo, a aposentadoria tende a ser muito baixa e pode não garantir continuidade para familiares.}

  Ele descreve que, na capitalização, ``se um deles fez o depósito e aos 85 anos o cidadão morreu, sabe qual o vencimento que ela fica? Zero!'', em referência à ausência de pensão. Também sustenta que, com contribuição reduzida, o trabalhador terminaria ``com exatamente um quarto, R\$750''.

  \item \textbf{Exemplos internacionais mostram que o sistema gerou problemas.}

  Ele diz que ``o sistema da capitalização já foi adotado no Chile, na Colômbia, no México e no Peru'' e que ``todos eles estão voltando atrás'', citando aposentadorias muito baixas e cobertura restrita. Como exemplo do Chile, menciona inclusive casos de ``suicídio entre os idosos'' e a ``média da aposentadoria no Chile'' como equivalente a ``meio salário mínimo''.

  \item \textbf{Ele usa dados e cenários para ilustrar o impacto sobre renda e o ganho do sistema financeiro.}

  No debate, cita uma estimativa segundo a qual ``a remuneração dos bancos consumiria R\$105.701,43'', o que equivaleria a mais de 62\% do patrimônio do empregado, deixando o trabalhador com menos. E conclui que ``um sistema sem empregador e com instituição financeira faturando um percentual será um fracasso, é condenar o trabalhador à miséria''.

  \item \textbf{Afirma que não há país que tenha aplicado esse modelo ``com recurso só do trabalhador'' e tido êxito.}

  Ele desafia: ``Se me apresentarem um único país que aplicou a tal de capitalização com o recurso só do trabalhador e que deu certo, eu mudo aqui o meu discurso''.
\end{enumerate}

\textbf{Exemplos usados por Paulo Paim}

\begin{itemize}
  \item Chile: aposentadorias baixas, menção a suicídio entre idosos e retorno parcial atrás do modelo.
  \item Chile, Colômbia, México e Peru: países citados como evidência de insucesso ou revisão do sistema.
  \item Cálculo ilustrativo em reais: estimativa de remuneração bancária consumindo grande parte do patrimônio.
  \item Cenário de pensão: caso em que o cônjuge remanescente ficaria com ``zero''.
  \item Desafio empírico: pedido de um único caso internacional bem-sucedido com recurso apenas do trabalhador.
\end{itemize}

\textbf{Limitações apontadas pela própria resposta}

\begin{itemize}
  \item O contexto não traz uma lista única e consolidada dos argumentos; eles aparecem distribuídos em vários trechos.
  \item A conclusão aparece dispersa em passagens distintas, e não em uma formulação final única.
\end{itemize}

\textbf{Referências retornadas pelo sistema}

\begin{itemize}
  \item Paulo Paim (PT/RS), 21/05/2019, \url{https://www25.senado.leg.br/web/atividade/pronunciamentos/-/p/texto/453686}
  \item Paulo Paim (PT/RS), 20/05/2019, \url{https://www25.senado.leg.br/web/atividade/pronunciamentos/-/p/texto/453766}
  \item Paulo Paim (PT/RS), 08/02/2019, \url{https://www25.senado.leg.br/web/atividade/pronunciamentos/-/p/texto/450736}
  \item Paulo Paim (PT/RS), 19/02/2019, \url{https://www25.senado.leg.br/web/atividade/pronunciamentos/-/p/texto/450992}
  \item Paulo Paim (PT/RS), 07/02/2019, \url{https://www25.senado.leg.br/web/atividade/pronunciamentos/-/p/texto/450672}
\end{itemize}

\subsection{Caso 2: reconhecimento de limite em desambiguação de autoria}

\textbf{Pergunta}

Quais propostas são realmente de Confúcio Moura e quais aparecem em discursos de outros senadores, mas poderiam ser confundidas com as dele?

\textbf{Resposta}

No contexto fornecido, \textbf{não é possível identificar com segurança} quais propostas são realmente de Confúcio Moura e quais aparecem em discursos de outros senadores, mas poderiam ser confundidas com as dele, porque o material recuperado traz apenas trechos e resumos de discursos, sem uma base completa que permita comparar cada proposta e distinguir, com segurança, autoria, citação ou mera concordância.

Dito isso, a resposta ainda consegue apontar, com base apenas no contexto recuperado, algumas linhas de atuação associadas a Confúcio Moura:

\begin{itemize}
  \item educação como prioridade, inclusive com a formulação de que ``a saída é a educação'';
  \item manifestação favorável à perenização do Fundeb e à ampliação de recursos;
  \item defesa de maior qualidade da educação básica e posição contrária à indicação política para diretorias escolares e secretarias municipais;
  \item intenção de proposta legislativa para melhorar condições de trabalho de catadores e recicladores.
\end{itemize}

Ao mesmo tempo, a resposta sinaliza corretamente o que \textbf{não foi possível confirmar}:

\begin{itemize}
  \item o contexto não traz uma separação completa entre ``propostas de Confúcio Moura'' e ``propostas de outros senadores'';
  \item alguns trechos aparecem em meio a debates ou falas de plenário, sem permitir distinguir com segurança autoria, apoio, comentário ou simples referência;
  \item por isso, o sistema evita afirmar além do que o material recuperado realmente sustenta.
\end{itemize}

\textbf{Por que este caso é relevante}

Este exemplo é metodologicamente importante porque mostra um tipo de bom comportamento que nem sempre aparece em demonstrações de sistemas RAG: em vez de forçar uma resposta completa sob incerteza, o sistema reconhece insuficiência de contexto, reduz o escopo da afirmação e preserva cautela na atribuição de autoria. Esse padrão é especialmente valioso em acervos parlamentares, nos quais temas semelhantes, vocabulário político recorrente e apartes em plenário podem induzir confusão entre autores.

\textbf{Referências retornadas pelo sistema}

\begin{itemize}
  \item Confúcio Moura (MDB/RO), 14/03/2019, \url{https://www25.senado.leg.br/web/atividade/pronunciamentos/-/p/texto/451323}
  \item Confúcio Moura (MDB/RO), 06/02/2020, \url{https://www25.senado.leg.br/web/atividade/pronunciamentos/-/p/texto/463649}
  \item Confúcio Moura (MDB/RO), 27/09/2019, \url{https://www25.senado.leg.br/web/atividade/pronunciamentos/-/p/texto/458160}
  \item Confúcio Moura (MDB/RO), 20/10/2020, \url{https://www25.senado.leg.br/web/atividade/pronunciamentos/-/p/texto/472033}
  \item Confúcio Moura (MDB/RO), 29/06/2021, \url{https://www25.senado.leg.br/web/atividade/pronunciamentos/-/p/texto/480774}
  \item Paulo Paim (PT/RS), 27/09/2019, \url{https://www25.senado.leg.br/web/atividade/pronunciamentos/-/p/texto/458158}
\end{itemize}

\subsection{Caso 3: pergunta multifocal e consolidação legislativa sob forte restrição de contexto}

\textbf{Pergunta}

Reconstrua, com base nos discursos recuperados, um plano legislativo consolidado para reduzir desigualdade no Brasil, indicando autores, medidas, cronograma, orçamento e consenso entre senadores.

\textbf{Resposta}

A resposta começa com um movimento metodologicamente correto: afirma que \textbf{não foi possível identificar, no conjunto de discursos recuperados, um plano legislativo único e consolidado} com autores, medidas, cronograma, orçamento e consenso explicitamente descritos dessa forma. Em vez de inventar uma síntese inexistente, o sistema reorganiza apenas o que o contexto realmente permite inferir, em cinco blocos: objetivo consolidado, medidas legislativas identificadas, cronograma, orçamento e consenso entre senadores.

\textbf{O que a resposta consegue fazer bem}

\begin{itemize}
  \item identifica objetivos recorrentes, como redução de desigualdades sociais, fortalecimento de políticas para população vulnerável, maior participação política de mulheres e busca de isonomia na distribuição orçamentária;
  \item recupera medidas legislativas concretas, citando, por exemplo, o PL n.\ 1.951/2021, a defesa da PEC n.\ 32/2022 e debates sobre orçamento secreto, emendas de bancada e limites de financiamento eleitoral;
  \item separa com clareza o que o contexto permite reconstruir e o que \textbf{não foi possível consolidar} com segurança, especialmente em relação a cronograma único, orçamento total e consenso formal entre senadores;
  \item evita apresentar como fato consumado um ``plano legislativo'' que, nos discursos recuperados, aparece apenas como conjunto fragmentário de propostas, críticas e agendas parciais.
\end{itemize}

\textbf{Por que este caso é mais desafiador}

Esta pergunta tensiona o sistema em múltiplas frentes ao mesmo tempo:

\begin{itemize}
  \item exige síntese entre vários autores e temas;
  \item pressupõe integração entre discursos que não necessariamente foram produzidos como parte de um mesmo plano;
  \item pede elementos estruturados --- cronograma, orçamento e consenso --- que raramente aparecem completos em falas parlamentares isoladas;
  \item aumenta o risco de o modelo ``costurar'' artificialmente uma política pública única a partir de trechos apenas parcialmente compatíveis.
\end{itemize}

\textbf{Limitações observadas}

\begin{itemize}
  \item embora a resposta seja cautelosa, ela já opera num nível de abstração maior do que os dois casos anteriores;
  \item a organização em blocos pode transmitir ao leitor a impressão de maior unidade programática do que aquela efetivamente sustentada pelos discursos;
  \item este é um bom exemplo de situação em que o sistema se comporta melhor ao reconhecer fragmentação do contexto do que ao tentar entregar uma resposta ``fechada''.
\end{itemize}

\textbf{Por que este caso é relevante}

Este terceiro exemplo ajuda a demonstrar um limite importante da solução: mesmo quando o sistema evita alucinação dura, perguntas amplas, multifocais e orientadas à consolidação normativa exigem supervisão humana mais forte. Em acervos parlamentares, a existência de temas comuns entre discursos não equivale, por si só, à existência de um plano legislativo único, consensual e plenamente estruturado.

\textbf{Referências retornadas pelo sistema}

\begin{itemize}
  \item Fabiano Contarato (REDE/ES), 10/04/2019, \url{https://www25.senado.leg.br/web/atividade/pronunciamentos/-/p/texto/452433}
  \item Orlando Silva (PCdoB/SP), 16/12/2022, \url{https://www25.senado.leg.br/web/atividade/pronunciamentos/-/p/texto/495399}
  \item Paulo Paim (PT/RS), 15/07/2021, \url{https://www25.senado.leg.br/web/atividade/pronunciamentos/-/p/texto/481744}
  \item Lasier Martins (PODEMOS/RS), 14/12/2021, \url{https://www25.senado.leg.br/web/atividade/pronunciamentos/-/p/texto/486225}
  \item Alessandro Vieira (PSDB/SE), 13/12/2022, \url{https://www25.senado.leg.br/web/atividade/pronunciamentos/-/p/texto/494690}
  \item Renan Calheiros (MDB/AL), 16/12/2022, \url{https://www25.senado.leg.br/web/atividade/pronunciamentos/-/p/texto/495400}
  \item Paulo Paim (PT/RS), 26/08/2019, \url{https://www25.senado.leg.br/web/atividade/pronunciamentos/-/p/texto/458186}
  \item Marcelo Castro (MDB/PI), 10/07/2019, \url{https://www25.senado.leg.br/web/atividade/pronunciamentos/-/p/texto/455943}
  \item Rodrigo Pacheco (PSD/MG), 20/12/2021, \url{https://www25.senado.leg.br/web/atividade/pronunciamentos/-/p/texto/486289}
\end{itemize}

% \clearpage
\phantomsection
\refstepcounter{section}
\addcontentsline{toc}{section}{APÊNDICE \Alph{section} -- PROMPT DO FLUXO RAG}
\begin{center}
\textbf{APÊNDICE \Alph{section} -- PROMPT DO FLUXO RAG}
\end{center}
\label{sec:apendice-prompt-rag}

Este apêndice registra o prompt configurado no modelo RAG do Open WebUI\footnote{\url{https://github.com/fabriciosantana/mcdia/blob/main/05-iag/4-project/eval/prompts/rag_prompt.md}}, associado ao uso interativo do sistema na própria plataforma e à bateria automatizada. O prompt explicita o núcleo metodológico adotado no experimento: fidelidade ao contexto recuperado, controle de escopo, reconhecimento de limites, cautela no uso de conhecimento externo e orientação à auditabilidade.

\subsection{Prompt operacional do modelo RAG no Open WebUI}

\begin{Verbatim}[breaklines=true]
### Tarefa

Responda a pergunta do usuário usando prioritariamente o contexto recuperado abaixo.

### Regras de resposta

- Use o contexto como fonte principal e autoritativa.
- Nao invente informações que nao estejam claramente apoiadas no contexto.
- Se a resposta estiver apenas parcialmente no contexto, responda somente a parte sustentada e diga explicitamente o que nao foi encontrado.
- Se usar conhecimento externo, sinalize explicitamente com a expressão: "Fora do contexto fornecido:".
- Nao misture conhecimento externo com fatos do contexto sem marcar a diferença.
- Quando houver múltiplos trechos relevantes, sintetize apenas os pontos realmente necessários para responder.
- Prefira os trechos mais diretamente relacionados a pergunta.
- Se houver conflito, ambiguidade ou baixa qualidade no contexto, diga isso claramente.
- Nao peca esclarecimento ao usuário se ja for possível responder parcialmente com segurança.
- Responda no mesmo idioma do usuário.
- Seja fiel ao conteúdo recuperado, mesmo que a resposta fique incompleta.
- Nao use conhecimento externo, exceto se o usuário pedir explicitamente.
- Se a resposta nao estiver no contexto, diga: "Nao encontrei essa informação no contexto fornecido."
- Antes de responder, verifique se cada afirmação relevante esta apoiada por pelo menos uma citação do contexto.
- Se nao estiver, remova a afirmação ou marque explicitamente que nao foi encontrada no contexto.

### Citações

- Coloque a citação imediatamente apos a afirmação correspondente.
- Use apenas as citações disponíveis no contexto.
- Nao inclua tags XML na resposta final.
- Nao agrupe citações no fim; distribua-as ao longo da resposta.

### Estilo de saída

- Responda de forma clara, objetiva e direta.
- Evite floreios e explicações desnecessárias.
- Nao copie trechos longos literalmente; prefira paráfrase fiel.
- Se a pergunta pedir enumeração, use lista.
- Se a pergunta pedir síntese, use um paragrafo curto seguido de pontos principais, se necessário.

Ao final de toda resposta, gere uma seção final obrigatória com o título:

### Referencias:

Nessa seção, liste os discursos usados para construir a resposta no formato:
- Nome do Senador (PARTIDO/UF) | Data: AAAA-MM-DD | Link: URL

- Use apenas os metadados disponíveis no contexto.
- Nao invente informações ausentes.
- Nao repita referencias.
- Se nenhum metadado estiver disponível, escreva:
-- Nao foi possível identificar as referencias completas no contexto fornecido.

### Estrutura da resposta

Siga esta ordem, quando aplicável:
1. Resposta direta a pergunta.
2. Evidencias principais do contexto.
3. Limitações do que nao foi possível confirmar.
4. Referencias
- Nome do Senador (PARTIDO/UF) | Data: AAAA-MM-DD | Link: URL

<context>
{{CONTEXT}}
</context>
\end{Verbatim}

\end{document}
